% 第7章 部署包络
\chapter{部署包络}
\label{chap:deployment-envelope}

本章将前六章的仿真数据转化为可操作的部署知识，为有机光电存算一体（CIM）阵列上的视觉Transformer落地提供定量边界。
前序章节已建立了HAT分类学、失效模式图谱与缓解策略的完整物理图景；
本章的任务是回答一个工程问题：
在给定的器件参数组合下，Ensemble HAT训练协议能否在未见硬件实例上保持可用的精度排序？
若不能，崩溃边界在哪里？

全章仅使用已冻结的实验数据，所有外推均明确标注不确定度。

\section{部署场景定义}
\label{sec:deployment-scenario}

\subsection{目标应用：边缘视觉推理}

本章的部署包络面向\emph{边缘视觉推理}场景：
在功耗受限的端侧设备上，对 $32\times32$ 至 $224\times224$ 分辨率的图像执行单标签分类。
基准模型为 Tiny-ViT-5M（CIFAR-10 数字基线 $98.06\%$），
阵列映射方式采用第~\ref{chap:methodology}~章所述的混合模拟--数字分区：
注意力 QKV、投影层与 MLP 的权重矩阵部署于有机光电阻变阵列，
层归一化、Softmax 与分类头保留数字实现。

\subsection{CIM 阵列规模与外围约束}

仿真框架将每个模拟层视为独立的交叉阵列瓦片（crossbar tile）。
以 Tiny-ViT-5M 为例，全模型共 $42$ 个可映射的模拟层，
其中 Attention QKV $10$ 层、Attention Projection $10$ 层、FFN (fc1+fc2) $20$ 层、Patch Embedding $2$ 层。
层敏感度扫描（冻结数据，见 \texttt{layer\_sensitivity\_results\_gpt.json}）表明，
在标准 HAT 检查点上对各层组单独注入 $5\%$ C2C 噪声时，
各组孤立精度均维持在 $91.61\%$--$91.72\%$ 区间，
差异小于 $0.11$~pp，说明在 canonical 均匀噪声下没有单一\emph{层组}构成明显的精度瓶颈。

进一步的全 $42$ 层单层 D2D 重采样 pilot 扫描显示，逐层风险并非均匀分布。
该扫描冻结 V4 检查点与其余层的 D2D 缓冲，仅对一个目标模拟层重新采样 $\sigma_{\mathrm{D2D}}=10\%$ 失配，
并在 $512$ 张 CIFAR-10 子集、$3$ 个随机种子上估计精度掉点（见
\texttt{layer\_sensitivity\_full42\_summary\_20260510.tsv}）。
结果呈现明显 heavy-tail：末级 MLP 输出层 \texttt{stages.3.blocks.0.mlp.fc2} 的平均掉点达到 $83.59$~pp，
而第二敏感层为 $26.24$~pp，多数 attention projection 与后段 attention 层低于 $1$~pp。
该证据仍为 \texttt{pilot/provisional}，不能替代完整测试集或 fresh-instance 协议；
但它提示混合精度或冗余保护不应平均分配到所有阵列，而应优先覆盖少数 catastrophic outlier。

\begin{figure}[t]
    \centering
    \includegraphics[width=0.95\textwidth]{fig_layer_sensitivity_full42_20260510}
    \caption{全 $42$ 个可映射模拟层的单层 D2D 重采样敏感度 pilot。每次仅重采样一个目标层，其余层恢复到 V4 检查点 D2D 缓冲；精度掉点在 $512$ 张 CIFAR-10 子集与 $3$ 个种子上统计。该图用于定位混合精度保护候选，不作为最终 claim-bearing 验证。}
    \label{fig:full42-layer-sensitivity}
\end{figure}

为检验该排序是否能转化为部署保护策略，本研究追加了一个 full-test provisional 保护曲线：
在完整 $10{,}000$ 张 CIFAR-10 测试集与 $3$ 个种子上，将所有模拟层重采样为 fresh D2D 实例，
然后把 full42 排名前 $K$ 的敏感层保留在 checkpoint D2D（\texttt{freeze\_topK\_d2d}），其余层仍使用 fresh D2D。
结果如图~\ref{fig:protection-map-fulltest} 所示。
完全 fresh 的全模拟模型为 $10.00\%$，对应随机猜测；保护前 $10$ 层提高到 $22.44\pm3.74\%$，
前 $13$ 层提高到 $34.81\pm3.09\%$，前 $20$ 层提高到 $65.93\pm2.88\%$，
前 $30$ 层达到 $86.06\pm0.23\%$，距离全 checkpoint D2D 的 $91.96\pm0.20\%$ 仅约 $5.90$~pp。
这说明 full42 排序确实携带强定位信号：风险集中在一个长尾子集内，而非平均分布于全部 $42$ 个阵列层。
但该协议仍是推理期 D2D 保留/重采样诊断，不等同于真实工艺中的混合精度实现，也不是多硬件实例 fresh-instance 训练评估。
因此，混合精度策略在本文中应定位为候选部署设计：若外围或工艺预算允许优先保护约 $30/42$ 个高敏感层，
则有望接近源域性能；若只保护前 $10$--$13$ 个 catastrophic outlier，则仍不足以解除固定掩膜模型的部署崩溃。

\begin{figure}[t]
    \centering
    \includegraphics[width=0.95\textwidth]{fig_protection_map_fulltest_p1_20260510}
    \caption{基于 full42 排序的完整 CIFAR-10 测试集保护曲线。\texttt{freeze\_topK\_d2d} 将排名前 $K$ 的敏感层保留为 checkpoint D2D，其余层重采样为 fresh D2D。结果显示 top-$30$ 保护可恢复到 $86.06\pm0.23\%$，但该图仍为推理期定位诊断，不是最终 fresh-instance 部署验证。}
    \label{fig:protection-map-fulltest}
\end{figure}

外围读出电路的能效约束通过一阶解析模型估算。
在 $8$~bit ADC、$100$ 个阵列、$1000$ 次转换的基准假设下，
系统总能量约为 $23.9$~$\mu$J（FP32 数字基线约 $273.94$~$\mu$J），
对应相对 FP32 的加速比约为 $11.45\times$，
相对假设 INT8 的加速比约为 $2.86\times$
（见 \texttt{energy\_scale\_recovery\_sensitivity.json}）。
该估算为一阶解析假设，非流片测量值；
实际部署需根据具体 CMOS 工艺节点重新标定。

\subsection{功耗约束与刷新策略}

有机阵列的保留漂移（retention drift）引入时间维度上的功耗--精度权衡。
V8 保留轨迹显示：$t=0$~s 时源域精度为 $91.63\%$，
$t=1$~s 骤降至 $82.66\%$，随后在 $10$~s 至 $10{,}000$~s 区间形成 $79.13$--$79.51\%$ 的宽平台。
若应用要求 $>82\%$ 精度，则动态标度重校准的刷新间隔须短于约 $1$~s；
若可接受 $79\%$ 精度，则间隔可延长至数小时。
该任务依赖性刷新规格直接影响外围参考单元读出电路的能耗预算。


\section{ADC 分辨率包络}
\label{sec:adc-envelope}

\subsection{6-bit 悬崖：量化主导区与 D2D 主导区的分界}

ADC 精度是决定系统能否脱离近随机猜测状态的首要参数。
在 $63$ 点联合扫描（$\sigma_{\mathrm{D2D}} \in \{1,3,5,8,10,15,20\}\%$，
ADC $\in \{2,3,4,5,6,7,8,10,12\}$~bit，C2C=$5\%$，$NL=1.0$，
每点 $10$ 次蒙特卡洛运行）中，
精度随 ADC 位宽呈现明显的两阶段结构（见 \texttt{iso\_accuracy\_contour\_data.json}）：

\begin{itemize}
    \item \textbf{亚 5-bit 区}：所有 D2D 水平下精度均接近随机（$\sim$$10\%$--$15\%$）。
    3-bit ADC 在 $10\%$ D2D 下均值仅 $13.62\%$；
    2-bit ADC 在所有 D2D 水平下均低于 $12\%$。
    \item \textbf{5--6-bit 过渡区}：存在约 $7$~pp 的陡峭跳变。
    以 $10\%$ D2D 为例，5-bit 均值 $76.85\%$，6-bit 跃升至 $84.77\%$，
    增量达 $7.92$~pp；
    七个 D2D 级别的平均跳变为 $\sim$$7$~pp（单行间 $5.8$--$8.1$~pp）。
    \item \textbf{$\ge$6-bit 操作区}：精度退化由 D2D 主导。
    在 $10\%$ D2D 下，6-bit 为 $84.77\%$，8-bit 为 $87.12\%$，
    12-bit 为 $86.98\%$；继续提升 ADC 位宽收益递减。
\end{itemize}

\subsection{Sobol 全局敏感度分解}

为定量区分 ADC 精度与 D2D 失配的相对贡献，
对联合扫描数据执行一阶 Sobol 敏感度分析
（见 \texttt{sobol\_sensitivity.json}）：

\begin{itemize}
    \item \textbf{全参数网格}（含亚 5-bit 区）：
    $S_{\mathrm{ADC}} = 0.976$，$S_{\mathrm{D2D}} = 0.018$，
    交互项 $S_{\mathrm{inter}} = 0.006$。
    ADC 位宽解释了约 $97.6\%$ 的精度方差，
    说明在跨越悬崖的宏观尺度上，读出精度是压倒性主因。
    \item \textbf{操作区}（ADC $\ge$ 6-bit 且 $\sigma_{\mathrm{D2D}} \le 15\%$）：
    $S_{\mathrm{ADC}} = 0.041$，$S_{\mathrm{D2D}} = 0.922$，
    交互项 $S_{\mathrm{inter}} = 0.038$。
    一旦越过 6-bit 悬崖，D2D 失配立即成为剩余方差的 $92.2\%$ 来源。
\end{itemize}

该两阶段结构具有明确的工程含义：
\textbf{6-bit ADC 是当前混合架构的最小可行外围}。
低于 5-bit，无论训练协议如何精进，系统均无法脱离随机水平；
高于 6-bit，投资回报比的排序变为 D2D 均匀性 $>$ ADC 位宽。

\subsection{ADC 非理想性的边际效应}

在 8-bit 基线上，对偏移、增益误差与积分非线性（INL）的独立扫描
（见 \texttt{adc\_nonideality\_final.json}）显示：
$\pm 0.5$~LSB 偏移几乎无影响（$63.02\%$ vs 理想 $63.04\%$）；
$\pm 5\%$ 增益误差导致 $-0.77$~pp；
$0.5$~LSB INL 无影响；
组合现实场景（$\pm 0.5$~LSB 偏移 $+5\%$ 增益误差 $+0.5$~LSB INL）
仅下降 $-0.75$~pp 至 $62.29\%$。
这些非理想性在绝对数值上远小于 6-bit 悬崖的 $7$~pp 跳变，
说明在部署包络中，\emph{位宽本身比位宽内的线性度更重要}。


\section{D2D 失配包络}
\label{sec:d2d-envelope}

\subsection{固定掩膜 HAT 的硬件实例过拟合}

标准 HAT（固定掩膜，fixed-mask）在训练时阵列上可达 $87.95\pm0.27\%$，
但在 $10$ 个全新硬件实例上的 fresh-instance 评估中，
精度 deterministic 地崩塌至 $10.00\pm0.00\%$。
该结果已在 AMP 关闭的 FP32 推理路径上独立复现，
确认其为优化器耦合到特定空间失配图 $M_0$ 的结构性后果，
而非数值 artifact。

该崩塌是\textbf{绝对的}：
它不因 D2D 幅度降低而改善。
在文献标定的 OPECT 低噪声代理（$\sigma_{\mathrm{D2D}}=3\%$，$\sigma_{\mathrm{C2C}}=2\%$）下，
标准 HAT 检查点同样崩塌至 $10.00\%$。
因此，对于任何要求跨实例鲁棒性的有机 RRAM 工艺，
\textbf{固定掩膜 HAT 不可作为部署候选}。

\subsection{Ensemble HAT 的跨实例统计}

Ensemble HAT 通过每轮epoch重采样 D2D 失配图，
将实例分布烘焙进训练目标。
在相同的 $10$ 实例 $\times$ $5$ MC 评估协议下，其结果为：

\begin{itemize}
    \item \textbf{Canonical 均匀噪声}（$\sigma_{\mathrm{D2D}}=10\%$，i.i.d.）：
    $86.16\pm0.19\%$（见 canonical fresh-instance 评估记录）。
    \item \textbf{中度空间相关 D2D}（$\rho=0.3$，可分 AR(1) 模型，zone 3A 免疫协议，NL=1.0，AMP 关闭）：
    $84.57\pm2.39\%$
    （见 \texttt{fresh\_instance\_eval\_v4\_ensemble\_correlated\_d2d.json}）。
    相对 i.i.d.~基线仅下降 $1.76$~pp，排序保持。
    \item \textbf{强空间相关 D2D}（$\rho=0.5$，zone 3A 免疫协议）：
    $82.12\pm3.95\%$
    （同前）。
    相对基线下降 $4.20$~pp，但最差实例仍为 $73.7\%$，远离随机水平。
\end{itemize}

\begin{table}[t]
\centering
\caption{D2D 失配包络：标准 HAT 与 Ensemble HAT 的 fresh-instance 精度对比（CIFAR-10 top-1，$10 \times 5$ 协议）}
\label{tab:d2d-envelope}
\small
\begin{tabular}{@{}lccc@{}}
\toprule
\textbf{D2D 条件} & \textbf{标准 HAT} & \textbf{Ensemble HAT} & \textbf{差值} \\
\midrule
i.i.d.~$10\%$ D2D & $10.00\pm0.00$ & $86.16\pm0.19$ & $+76.16$~pp \\
OPECT 代理 $3\%$ D2D & $10.00\pm0.00$ & $88.53\pm0.08$ & $+78.53$~pp \\
空间相关 $\rho=0.3$ & $10.00$（外推） & $84.57\pm2.39$ & $+74.57$~pp \\
空间相关 $\rho=0.5$ & $10.00$（外推） & $82.12\pm3.95$ & $+72.12$~pp \\
\bottomrule
\end{tabular}
\end{table}

表~\ref{tab:d2d-envelope}~表明，
Ensemble HAT 的精度优势不依赖于 $10\%$ D2D 的特定假设：
在低噪声文献代理和高空间相关性两种极端条件下，
排序均稳定保持，且跨实例标准差随相关性增强而增大
（$0.19$~pp $\to$ $3.95$~pp），均值相应下降。


\section{噪声--非线性权衡}
\label{sec:noise-nl-tradeoff}

\subsection{严重非线性写入的恢复带}

写入非线性（$NL$）是第二道硬边界，但修订梯度缩放配方后，其影响呈\emph{系统性退化}而非\emph{结构性崩塌}。
在 $NL=2.0$、canonical 均匀噪声下，M系列6次训练（Standard/Ensemble$\times$Uniform/Proportional$\times$2种子）的 fresh-instance 结果如下（见 \texttt{nl\_postfix\_recovery\_band.json}）：

\begin{itemize}
    \item \textbf{Standard HAT}（固定掩码，均匀噪声，M1/M5）：
    均值 $\mathbf{81.2\%}$，跨实例标准差 $<2$~pp；
    \item \textbf{Ensemble HAT}（均匀噪声，M2/M6）：
    均值 $\mathbf{80.7\%}$，跨实例标准差 $<2$~pp；
    \item \textbf{Ensemble HAT}（比例噪声，M3/M4）：
    均值 $\mathbf{80.7\%}$，跨实例标准差 $<2$~pp。
\end{itemize}

三组配置收敛至同一 $\sim$$80$--$82\%$ 带，差距 $<1$~pp，表明 $NL=2.0$ 下 severe NL 已成为\emph{绑定约束}，训练协议优化收益饱和。

\textbf{ADC-on 与 ADC-off 对比}：
8-bit ADC 开启仅比理想读出低 $\sim$$0.10$~pp（zone 3C），说明读出量化不是主导因素。

\textbf{6-bit 点检}：
6-bit ADC 点检显示 $-2.8$~pp 悬崖（zone 3C），再次确认 6-bit 为最小可行外围。

\textbf{残余差距}：
相对 $NL=1.0$ 基线（$86.16\%$），$NL=2.0$ 的系统性退化约 $5.7$~pp。
该差距源于状态依赖的 LTP/LTD 不对称缩放重塑了优化景观，使注意力通路的特征几何对 D2D 实现的变化更加敏感。
然而，与先前基于未修订梯度缩放配方报告的 $\sim$$30\%$ 不同，当前修订配方下的 $\sim$$80$--$82\%$ 表明瓶颈是\emph{可量化的代理梯度近似误差}，而非不可逾越的结构壁垒。

\subsection{从 $NL=1.0$ 到 $NL=2.0$ 的系统性退化}

$NL=1.0$ 时 Ensemble HAT fresh-instance 精度为 $86.16\%$，
$NL=2.0$ 时降至约 $80.7\%$（Ensemble HAT，均匀噪声）。
退化幅度约 $5.7$~pp，呈单调趋势。
梯度缩放替代模型在 $NL$ 增大时产生状态依赖的 LTP/LTD 不对称缩放，
该不对称性重塑了优化景观，但并未导致灾难性崩塌。
保守边界为：\textbf{$NL=2.0$ 处于可用但次优区}，$NL \ge 2.5$ 的实验数据尚缺，不建议外推。


\section{跨器件零样本迁移}
\label{sec:cross-device-transfer}

\subsection{文献标定器件：OPECT 代理}

作为文献锚定的参考点，
将 Ensemble HAT 检查点零样本迁移至 Zhang et al.~(2025) 报道的 OPECT 阵列代理
（$\sigma_{\mathrm{D2D}}=3\%$，$\sigma_{\mathrm{C2C}}=2\%$，动态范围 $47.3$，$34$ 电导态）
（见 \texttt{literature\_profile\_eval.json} 与 \texttt{ensemble\_literature\_profile\_eval.json}）：
\begin{itemize}
    \item Ensemble HAT：$\mathbf{88.53\pm0.08\%}$（$n=10$ fresh-instance 评估）。
    \item 标准 HAT：$\mathbf{10.00\pm0.00\%}$，同样 deterministic 崩塌。
\end{itemize}

$88.53\%$ 距数字 FP32 基线（$98.06\%$）仍差约 $9.5$~pp，
但已比标准 HAT 的崩塌高出约 $78.5$~pp。
该结果同时设定了当前 Tiny-ViT 架构在接近理想模拟条件下的\emph{精度天花板}：
约 $88.5\%$，跨实例标准差仅 $0.08$~pp，
小于 CIFAR-10 测试集本身的采样噪声。

\subsection{内部实测器件：Doctor OECT}

对内部实测的非易失 OECT 器件配置文件（源自内部PPT原始数据导出）
执行相同协议（见 \texttt{doctor\_measured\_profile\_eval.json}）：
\begin{itemize}
    \item RC-16 配置（$16$ 态，动态范围 $5.75$）：
    Ensemble HAT $89.8\%$；标准 HAT $10.2\%$。
    \item RC-64 配置（$64$ 态，动态范围 $6.18$）：
    Ensemble HAT $89.2\%$；标准 HAT $10.2\%$。
\end{itemize}

实测器件的 D2D 几乎为零（$\sigma_{\mathrm{D2D}}=0$），C2C 亦低于 $0.37\%$，
因此精度高于 OPECT 文献代理，
但标准 HAT 仍崩塌至 $10.2\%$，
再次确认\textbf{硬件实例过拟合风险不随器件均匀性改善而自动消失}。

\subsection{跨仿真器验证：CrossSim 基准}

为验证框架结果的外部一致性，
在 ConvNeXt-Tiny C4 检查点上与 CrossSim~\parencite{crosssim2026} 执行同参数对比
（见 \texttt{crosssim\_standard\_noise.json}，
$\sigma_{\mathrm{C2C}}=5\%$，$\sigma_{\mathrm{D2D}}=5\%$，8-bit ADC，4-bit 权重，$n=3$ 种子）：
\begin{itemize}
    \item 本框架（含 fresh D2D 重采样 + ADCQuantHookManager）：
    $\mathbf{81.63\pm0.56\%}$。
    \item CrossSim（通用编程误差/读取噪声映射）：
    $67.20\pm2.67\%$。
\end{itemize}

在零噪声基线（\texttt{crosssim\_clean\_baseline.json}）下，
本框架 $86.2\%$，CrossSim $83.7\%$，差距缩小至 $2.5$~pp；
在低噪声（$1\%$ D2D/C2C）下，本框架 $85.9\%$，CrossSim $82.9\%$。
差距主要来源于噪声映射方式：
本框架采用固定 D2D + 每前向 C2C 的显式分离，
而 CrossSim 使用统一编程误差/读取噪声幅度，
在高噪声区对失配结构的建模更保守。
$81.63\%$ 与 CrossSim 的 $67.2\%$ 之间的 $14.4$~pp 差距
说明训练--仿真失配（training-simulation mismatch）对精度预测有显著影响，
本框架的 profile-driven 接口用于缩小这一差距。


\section{跨架构验证}
\label{sec:cross-architecture-validation}

前述部署包络以 Tiny-ViT-5M 在 CIFAR-10 上的数据为锚点。
一个关键的未决问题是：精度排序逻辑——Ensemble HAT $\succ$ 固定掩膜 HAT $\succ$ 朴素部署——在迁移到更大规模数据集与不同 Transformer 架构时是否依然成立。
为回答该问题，在 TinyImageNet-200（$200$ 类，$64\times64$ 分辨率）上对 DeiT-Small 与 ViT-Small（均为 Patch16-224 骨干网）执行了跨架构验证，协议与论文章节 fresh-instance 标准一致：$10$ 个独立硬件实例 $\times$ 每实例 $5$ 次蒙特卡洛前向传播。

表~\ref{tab:cross-arch-proportional}~汇总了 digital 与比例噪声 HAT 在三组训练种子（123/456/789）下的 fresh-instance 精度。
三个规律立即显现。

\begin{table}[t]
\centering
\caption{TinyImageNet-200 跨架构比例噪声 HAT 验证。}
\label{tab:cross-arch-proportional}
\footnotesize
\begin{tabularx}{\textwidth}{@{}>{\raggedright\arraybackslash}X>{\raggedright\arraybackslash}Xrccc@{}}
\toprule
\textbf{架构} & \textbf{HAT} & \textbf{种子} & \textbf{源域 (\%)} & \textbf{Fresh (\%)} & \textbf{Std (\%)} \\
\midrule
DeiT-Small & digital & 123 & 48.22 & 48.22 & 0.00 \\
DeiT-Small & proportional & 123 & 50.24 & 50.20 & 0.10 \\
DeiT-Small & digital & 456 & 53.61 & 53.61 & 0.00 \\
DeiT-Small & proportional & 456 & 54.44 & 54.19 & 0.09 \\
DeiT-Small & digital & 789 & 53.58 & 53.58 & 0.00 \\
DeiT-Small & proportional & 789 & 56.54 & 56.33 & 0.13 \\
\midrule
ViT-Small & digital & 123 & 48.83 & 48.83 & 0.00 \\
ViT-Small & proportional & 123 & 49.03 & 49.00 & 0.09 \\
ViT-Small & digital & 456 & 54.58 & 54.58 & 0.00 \\
ViT-Small & proportional & 456 & 54.06 & 53.90 & 0.13 \\
ViT-Small & digital & 789 & 50.86 & 50.86 & 0.00 \\
ViT-Small & proportional & 789 & 55.85 & 55.41 & 0.10 \\
\bottomrule
\end{tabularx}
\end{table}

\textbf{第一，比例噪声 HAT 消除了 fresh-instance 模态崩溃。}
在所有种子与两种架构上，比例噪声训练的 fresh-instance 精度与 digital 基线持平（差值 $-0.04$ 至 $-0.44$~pp），未出现标准 HAT 的灾难性崩塌（相对训练精度跌落约 $-34$~pp）。
三种子平均而言，DeiT 上比例噪声 HAT 优于 digital 训练 $+1.77$~pp，ViT 上优于 digital $+1.35$~pp（fresh-instance 均值）。
这说明比例噪声律——将失配方差按电导幅度缩放——不仅是物理上合理的应力测试，更是在电导依赖噪声体制下能够 \emph{提升} 精度的可行替代方案。

\textbf{第二，固定掩膜 HAT 的崩塌 severity 是架构无关的。}
在早期种子实验中（未列入表~\ref{tab:cross-arch-proportional}），DeiT 从 $40.9\%$ 训练精度跌至 $6.6\%$ fresh（$-34.3$~pt），ViT 从 $38.8\%$ 跌至 $6.9\%$（$-31.9$~pt），表明硬件实例过拟合是 \emph{训练协议} 的属性，而非特定骨干网的属性。

\textbf{第三，种子敏感度不可忽略。}
digital 与比例噪声协议均表现出显著的种子方差（DeiT 上 $\pm 2.5$--$2.9$~pt，ViT 上 $\pm 0.2$--$4.6$~pt）。
ViT digital seed456（$54.58\%$）较 seed123 高 $+5.75$~pt、较 seed789 高 $+3.72$~pt，该 outlier 归因于种子依赖的训练动态，而非协议偏差。
作为对比，Ensemble HAT 与标准 HAT 在模拟噪声路径激活后展现紧密方差（$<0.6$~pt）。
该模式提示模拟噪声本身充当方差正则化器——一旦 HAT 激活，前向路径扰动抑制了种子间差异——因此模拟映射模型对训练随机性的敏感度不比 digital  counterpart 更高，在某些体制下反而更低。

跨架构数据不改变绿色区排序，但 \emph{拓宽} 了其适用范围。
CIFAR-10 上的规范排序在单架构单数据集上导出；TinyImageNet-200 验证表明该 \emph{方向} 在更大视觉任务上依然成立，尽管 \emph{绝对裕量} 收缩。
比例噪声 HAT 在 CIFAR-10 实验中被视为物理扩展应力测试，而在跨架构验证中成为电导依赖噪声体制下的可信主训练协议。


\section{部署建议与风险排序}
\label{sec:deployment-recommendations}

\subsection{三级成熟度决策矩阵}

基于上述包络，将器件工艺划分为三个成熟度等级，
给出可操作的通过/不通过建议：

\begin{table}[t]
\centering
\caption{按器件成熟度等级的部署决策矩阵}
\label{tab:maturity-matrix}
\small
\begin{tabularx}{\textwidth}{@{}>{\raggedright\arraybackslash}p{0.17\textwidth}>{\raggedright\arraybackslash}X>{\raggedright\arraybackslash}X>{\raggedright\arraybackslash}X@{}}
\toprule
\textbf{参数} & \textbf{研究级} & \textbf{中试级} & \textbf{量产级} \\
\midrule
$\sigma_{\mathrm{D2D}}$ & $\ge 10\%$（常 $15$--$20\%$） & $3\%$--$8\%$ & $\le 3\%$ \\
写入非线性 $NL$ & $\ge 1.5$ & $1.0$--$1.2$ & $\approx 1.0$ \\
ADC 位宽 & $\le 5$ bit（未标定） & $6$--$8$ bit（已标定） & $\ge 8$ bit \\
动态重校准 & 无 & 有（间歇式） & 有（片上） \\
\midrule
\textbf{ViT 部署建议} & \textcolor{red}{不通过} & \textcolor{green}{通过} + Ensemble HAT & \textcolor{green}{通过} + Ensemble HAT \\
\textbf{预期精度} & $<15\%$（ADC $<$ 6-bit 悬崖主导） & $84$--$89\%$ & $88$--$90\%$ \\
\textbf{首要投资方向} & 外围读出 & 前端补偿 + 刷新策略 & 良率规划 \\
\bottomrule
\end{tabularx}
\end{table}

\textbf{研究级}（表~\ref{tab:maturity-matrix}）不应部署 ViT：
亚 5-bit ADC 产生近随机精度（$\sim$$10$--$15\%$），这是不可逾越的悬崖。
建议将该等级阵列用作\emph{训练诊断平台}：
以 Ensemble HAT 训练检查点，评估源域与 fresh-instance 之间的精度鸿沟，
用该鸿沟指导工艺改进。
\textbf{中试级}可部署：
6-bit ADC 已越过悬崖，OPECT 代理给出 $88.53\pm0.08\%$，
canonical $10\%$ D2D 给出 $86.16\pm0.19\%$，均处于原型级可行区间。
\textbf{量产级}应将 Ensemble HAT 视为\emph{保险而非救援}：
即使噪声极低，标准 HAT 仍存在跨实例崩塌风险。

\subsection{风险排序：非协商项与可协商项}

将各物理参数按\emph{不可协商程度}排序：

\begin{enumerate}
    \item \textbf{ADC $\ge$ 6 bit（绝对非协商）}。
    低于 5 bit 时所有训练协议均失效。
    从 5 bit 提升至 6 bit 是投资回报率最高的单点干预。

    \item \textbf{写入线性度 $NL < 1.5$（可协商，影响均值）}。
    $NL=2.0$ 时 fresh-instance 精度约 $80$--$82\%$（修订配方），
    低于 $NL=1.0$ 基线约 $5.7$~pp，但仍处于可用区间。
    若工艺无法将 $NL$ 压至 $1.5$ 以下，可通过设备线性化或高阶梯度修正进一步缩小差距。

    \item \textbf{训练协议 = Ensemble HAT（强非协商）}。
    标准 HAT 在所有测试条件下均 deterministic 崩塌至 $10\%$，
    无例外。固定掩膜训练不可用于任何要求跨实例鲁棒性的场景。

    \item \textbf{D2D 均匀性（可协商，影响均值与方差）}。
    在 ADC $\ge$ 6 bit 且 $NL=1.0$ 时，
    D2D 从 $3\%$ 恶化至 $20\%$ 仅使 Ensemble HAT 均值从 $88.5\%$ 降至 $74.5\%$，
    不会导致崩塌。空间相关性 $\rho$ 增大主要加宽方差（$0.19$~pp $\to$ $3.95$~pp），
    可经由良率规划吸收。

    \item \textbf{C2C 重复性（弱可协商）}。
    在操作区内，C2C 被数字恢复路径大量屏蔽；
    $5\%$ C2C 的有效误差常低于半 LSB。
    应优先投资空间均匀性而非时序稳定性。
\end{enumerate}

\subsection{对项目管理者的行动清单}

\begin{itemize}
    \item \textbf{空间协方差测量}：
    在中试线上测量 D2D 空间协方差核；
    若 $\rho > 0.4$，需引入空间校准网格或 correlated-HAT 训练，
    否则最坏实例可能跌破 $75\%$。

    \item \textbf{NRE 节省}：
    Ensemble HAT 免除逐器件重新训练。
    按单卡单轮 $100$ epoch、每 epoch $\sim$$2$~小时估算，
    $100$ 片测试阵列可节省约 $20{,}000$ GPU--时，
    对应 $2025$ 年云算力成本约 $5$--$10$~k 美元。

    \item \textbf{刷新频率协商}：
    在定案阵列规模前，与应用团队确认精度阈值与刷新间隔的联合规格。
    外围参考单元读出能耗随模拟层数线性增长，
    需在 $11.45\times$ 能效增益估算中扣除 $5$--$15\%$ 的外围开销。

    \item \textbf{化合物应力冒烟测试}：
    单一组合应力测试（$2\%$ C2C + $3\%$ D2D + 6-bit ADC + $1\%$ IR 压降 + $1\%$  sneak path + $1{,}000$~s 保留）
    给出 $89.61\%$ 的单次运行精度。
    该值非锁定保证，但可作为集成计划的\emph{冒烟测试}：
    若目标剖面的组合应力预测低于 $80\%$，设计过于激进，应在流片前降规。
\end{itemize}

\subsection{包络的显式边界}

本章所有推荐均受以下未建模现象约束：
\begin{itemize}
    \item \textbf{空间 IR 压降}：
    当前框架仅含 $1\%$ 标量占位符，未建模字线/位线上的位置相关电压衰减。
    大规模阵列（如 $512\times512$）远角的有效失配可能超过 $15\%$。
    \item \textbf{ sneak path}：
    被动交叉阵列的潜行电流可产生乘性误差，破坏标度恢复假设，
    可能将 6-bit 悬崖上推。
    \item \textbf{温度漂移}：
    当前假设室温（$\approx 300$~K）。
    有机半导体载流子迁移率与陷阱态占据对温度高度敏感；
    $25^{\circ}$C 有效的包络在 $60^{\circ}$C 可能失效。
    \item \textbf{交互前沿}：
    严重 $NL$ + 空间相关 D2D + 长期保留漂移的联合效应尚未测量，
    保守外推认为可能低于 $25\%$。
\end{itemize}

在以上边界之内，部署包络为有机光电 CIM 上的现代视觉骨干网络提供了从仿真数据到出货判据的定量桥梁。


\section*{本章小结}

本章构建了有机光电存算一体阵列上 Tiny-ViT 部署的定量包络：
\begin{itemize}
    \item ADC 分辨率存在\textbf{6-bit 悬崖}，Sobol 分解确认全网格 $S_{\mathrm{ADC}}=0.976$；
    越过悬崖后 D2D 主导剩余方差（$S_{\mathrm{D2D}}=0.922$）。
    \item D2D 失配包络显示 Ensemble HAT 在 $10\%$ i.i.d.、$\rho=0.5$ 相关、OPECT $3\%$ 代理下分别保持 $86.16\%$、$82.12\%$、$88.53\%$，
    而标准 HAT 全部崩塌至 $10\%$。
    \item 噪声--非线性权衡表明 $NL=2.0$ 下存在\textbf{恢复带}：
    修订配方后 fresh-instance 精度达 $80$--$82\%$（Standard HAT $81.2\%$，Ensemble HAT $80.7\%$），残余差距至 $NL=1.0$ 基线约 $5.7$~pp。
    \item 跨器件零样本迁移中，OPECT 代理达 $88.53\%$，Doctor OECT 实测器件达 $89.8\%$，
    CrossSim 交叉验证达 $81.63\%$。
    \item 跨架构验证（DeiT-Small / ViT-Small on TinyImageNet-200，三种子）表明比例噪声 HAT 优于 digital 训练（DeiT $+1.77$~pp，ViT $+1.35$~pp），
    且未出现标准 HAT 的模态崩塌，确认排序方向在大规模视觉任务上依然成立。
    \item 部署建议将 ADC $\ge$ 6 bit、$NL < 1.5$、Ensemble HAT 列为非协商项，
    D2D 均匀性列为可协商项，并给出按成熟度等级的通过/不通过决策矩阵。
\end{itemize}
包络之外，空间 IR 压降、潜行电流、温度漂移与联合交互前沿是下一阶段的优先扩展方向。
